\chapter{Závěr}
Modelářské polohové servo spojuje motor, převodovku, snímač a regulaci v~jednom pouzdře. K~řízení mu stačí napájení a pulzy s~odpovídající délkou. Průmyslový servopohon má zpravidla samostatný zesilovač a jiné rozhraní; kontinuální RC servo zase neumožňuje bez doplněného měření přímo zadat úhel. Převodovka není samostatnou třídou použití, ale konstrukční vlastností obou skupin pohonů.

Pro lehký držák jsme zvolili polohové provedení MG996R 180°. Oproti menším SG90 a MG90S nabízí větší momentovou rezervu, vyžaduje však silnější napájení. Výslednou přesnost budou ovlivňovat vůle převodu, zatížení a kalibrace.

Samostatný generátor s~jedním NE555 a diodou nastavuje délku povelu potenciometrem; současně se mírně mění opakovací perioda. Je vhodný pro ruční tester a názorné vysvětlení řízení. Jeho omezením jsou tolerance součástek, nutnost kalibrace a chybějící programové omezení pohybu. Arduino Uno R3 nabízí stejné rozhraní pro servo, umožňuje nastavit meze v~programu a později doplnit filtraci nebo omezení rychlosti povelu. Pro výsledný přípravek je proto vhodnější.

Významnou část návrhu tvoří napájení a zemní návraty. Servo a řízení využívají společný 5V zdroj se samostatnými přívody; společná zem musí poskytovat referenci signálu, aniž by proud motoru procházel řídicí deskou. Zemní smyčku přes další přístroje je nutné rozlišit od úbytku na společném vodiči. Blokovací kondenzátory, krátké vedení a správné uspořádání doplňují dostatečný zdroj, ale nenahrazují jej.

Navržené řešení lze použít pro výukový přípravek, nastavování senzoru nebo lehký modelový mechanismus. Práce představuje návrh obvodu a programu; fyzická realizace nebyla provedena.
